\section{Methods}
\label{sec:methods}

The same preprocessing functions are called by the dataset builder, the training and evaluation code and the web service, so training and serving cannot drift apart. Sections~\ref{sec:preproc}--\ref{sec:exttask} describe the classifier and its evaluation; Sections~\ref{sec:probe}--\ref{sec:eval-order} describe the source-information analyses.

\subsection{Preprocessing}
\label{sec:preproc}
Each image is decoded to an 8-bit RGB array and passes through four deterministic steps. (i) A $3\times3$ Gaussian pre-blur with $\sigma_1=0.8$ suppresses noise before enhancement. (ii) Contrast-limited adaptive histogram equalisation (CLAHE)~\cite{pizer1987adaptive,zuiderveld1994contrast} with clip limit 2.0 and $8\times8$ tiles is applied to the lightness channel of the CIE $L^*a^*b^*$ representation. (iii) A Gaussian low-pass filter with $\sigma_2=1.0$ attenuates spatial frequency $(u,v)$, in cycles per pixel, by
\begin{equation}
H(u,v)=\exp\!\left(-2\pi^2\sigma_2^2\,(u^2+v^2)\right).
\label{eq:lowpass}
\end{equation}
The value $\sigma_2=1.0$ was chosen in the earlier V5 work on V5 validation data and was not re-tuned on V6. (iv) The image is resized to $224\times224$ with Lanczos interpolation and normalised with ImageNet statistics. The three filters run at native resolution before resizing, so images of different resolution receive different effective filtering; the $224\times224$ result is cached per image, keyed by the SHA-256 of the raw file. Augmentation during training only is a horizontal flip and a rotation within $\pm10^\circ$.

\subsection{Network, training and checkpoint selection}
The classifier is a DenseNet-121~\cite{huang2017densely} with ImageNet weights~\cite{deng2009imagenet}. With $\mathbf{z}\in\mathbb{R}^{1024}$ the global-average-pooled (GAP) backbone feature, the head computes
\begin{equation}
\mathbf{a}=\mathrm{Dropout}_{0.3}\!\left(\mathrm{ReLU}\!\left(W_1\,\mathrm{BN}(\mathbf{z})+\mathbf{b}_1\right)\right),\qquad
\mathbf{p}=\mathrm{softmax}(W_2\mathbf{a}+\mathbf{b}_2),
\label{eq:head}
\end{equation}
with $W_1\in\mathbb{R}^{256\times1024}$ and $W_2\in\mathbb{R}^{6\times256}$. The softmax scores are not calibrated probabilities of disease. The loss is the class-weighted cross-entropy with balanced weights $w_k=n/(K n_k)$. Training uses Adam~\cite{kingma2015adam}, mini-batches of 32 and seed \RSeed. Phase~1 (two epochs) trains the head with learning rate $10^{-3}$ and the backbone frozen; phase~2 unfreezes the last 30 backbone layers at learning rate $10^{-5}$ until three epochs pass without an improvement in validation macro-F1, an epoch limit, or a wall-clock budget. After every epoch the model is evaluated on the validation split, and the epoch with the highest validation macro-F1 (macro ROC-AUC as tie-breaker) is kept. Neither the test split nor the external cohort is used for training or checkpoint selection.

\subsection{Evaluation metrics}
\label{sec:metrics}
We report accuracy, per-class precision, recall and F1, macro-F1, one-vs-rest ROC-AUC (macro-averaged) and the expected calibration error (ECE) with 15 equal-width bins~\cite{naeini2015obtaining,guo2017calibration}. Proportions carry Wilson 95\% intervals~\cite{wilson1927probable}; for composite metrics we use percentile bootstrap intervals~\cite{efron1994introduction}. Both describe sampling variation of the evaluated set for a fixed model, not variation due to training.

\subsection{External binary task}
\label{sec:exttask}
Montgomery contains only normal and TB radiographs. A prediction is \emph{abnormal} if it is not Normal. We report sensitivity (TB scans predicted abnormal) always together with specificity (normal scans predicted Normal), the exact Tuberculosis recall, the ROC-AUC of $1-p_{\text{Normal}}$ and the full table of predicted classes. No threshold was tuned on the external cohort or the internal test split. Grad-CAM~\cite{selvaraju2017gradcam} maps were inspected qualitatively for one test image per class, chosen by a rule fixed before any heatmap was viewed (the correctly classified image with the median top-class score); saliency maps on chest radiographs are themselves unreliable~\cite{arun2021assessing}, so no conclusion is drawn from them beyond illustration.

\subsection{Source-identifiability probe}
\label{sec:probe}
To test whether the source of an image is recoverable from the features a classifier uses, we drew, per (source, class) pair, at most 150 training-split and 100 test-split images (fixed seed), giving \RProbeNTrain{} probe-training and \RProbeNTest{} probe-test images. Each cached image was mapped to the 1024-dimensional GAP vector by two frozen extractors: DenseNet-121 with ImageNet weights, which has learned nothing from this project, and the backbone of the selected model. After standardisation, a multinomial logistic regression ($C=1$, class weights inversely proportional to source frequency) was fitted to predict the source. We report balanced accuracy on the probe-test images, overall and within each class that has at least two sources. A high value means that source information is available in the features; it does not show that the trained classifier uses it, so the probe is indirect evidence.

\subsection{C1: spectral map of source and disease information}
\label{sec:meth-c1}
\paragraph{Bands} The cached $224\times224$ images were decomposed in the Fourier domain into six octave bands A--F by radial frequency $r$ (cycles per pixel) with band edges at $1/64$, $1/32$, $1/16$, $1/8$ and $1/4$ (A: below $1/64$, including the lowest frequencies; F: above $1/4$, with periods of 4 to 2 pixels). Edges have smooth raised-cosine transitions of $\pm0.25$ octave, so the six masks sum to one at every frequency. The DC coefficient is kept in every variant. For each band, a \emph{band-pass} image keeps only that band and a \emph{band-stop} image removes only that band. Band-pass images are not rescaled in contrast.

\paragraph{Features, probes and scores} Frozen 1024-dimensional GAP features were extracted from the full image and from the 12 band variants. The primary extractor is ImageNet DenseNet-121. Secondary checks are the backbone of the selected model (``trained backbone''; it saw all six sources in training, so it is not a held-out test) and ImageNet features of the no-filter cache of run~5 (the cached images without the three filters of Section~\ref{sec:preproc}; a sensitivity check). Two logistic-regression probes (standardised features, $C=1$, balanced class weights) were fitted on a fixed training sample of \RSpecMainTestNTrain{} images per variant, one predicting the source and one the disease, and scored on a validation or test sample. The test sample has \RSpecMainTestN{} images. We report the source balanced accuracy (chance $1/6=16.7\%$) and the primary disease score, the \emph{within-source} macro-F1: macro-F1 over the qualifying classes of each source, averaged over \RSpecMainTestNWithinSources{} sources (a (source, class) cell qualifies with at least 20 images). Because disease and source are confounded, the six-class macro-F1 can be bought with source cues and is secondary. The within-Normal source accuracy, the accuracy of the source probe on Normal images only, shows whether source information survives when disease is held fixed. Band-stop differences are $dS=\text{BA}_{\text{src}}(\text{full})-\text{BA}_{\text{src}}(\text{stop})$ and $dD$ for the within-source macro-F1 (positive means information lost). Intervals are percentile bootstrap 95\% intervals over the scored images (1,000 resamples) with one probe fit per variant; they therefore describe test-sample variation, not training variance.

\paragraph{Caveat on resampling} The cache is resampled to $224\times224$ after filtering, and sources differ in native resolution and field of view. The band map is therefore a map of the cached images, and part of the band structure may reflect resampling and pre-filters rather than acquisition physics.

\subsection{C3: linear source erasure on frozen features}
\label{sec:meth-c3}
We applied INLP~\cite{ravfogel2020null} to the standardised 1024-dimensional features of the full image (training, validation, test and Montgomery images). At each iteration a logistic regression predicting the six sources is fitted on the current training features, the row space of its coefficients is removed by orthogonal projection, and the step is repeated; the cap is 20 iterations, with early stopping if a fresh source classifier reaches validation balanced accuracy of at most $1/6+0.05$. After each iteration a new disease head (multinomial logistic regression) is retrained on the projected features, and the source accuracy a fresh probe can reach is recomputed. As a one-row check we applied closed-form LEACE~\cite{belrose2023leace} with the source labels. Scores are computed on the full validation and test splits; changes in the disease scores carry paired bootstrap intervals over test images. Montgomery was scored once, with the heads at $k=0$ and at the final iteration $k^\ast$. INLP is not claimed as a method of ours: it is applied as a measurement of how entangled source and disease information are.

\subsection{C4: linear-head leave-one-source-out}
\label{sec:meth-c4}
Following the leave-one-domain-out design of Yao et al.~\cite{yao2019strong}, a multinomial logistic regression on standardised frozen features was trained on the training splits of five sources and scored on the test split of the sixth, for each source in turn (``linear-head leave-one-source-out on frozen features''). The regularisation $C$ was chosen on the validation images of the other sources. ImageNet features are primary; the trained-backbone matrix is secondary and is never called held-out or external. A source is a held-out row if it has at least two evaluated classes, each with at least 20 test images in that source and 50 training images in the others; \RLosoImgNIncluded{} sources qualify (the COVID-19 Radiography Database has one class). Macro-F1 is computed on the evaluated classes only, and a prediction of any other class counts as an error. The in-distribution reference is the same recipe trained on all six sources and scored on the same images; the \emph{gap} is reference minus held-out. A binary abnormal-versus-normal view uses $1-P(\text{Normal})$. Intervals are bootstrap intervals over held-out test images. Montgomery was scored once with the head trained on all six sources.

\subsection{C2-lite and the decision rule fixed in advance}
\label{sec:meth-c2}
A full C2 (band-stop of the most source-heavy band, then retraining) was conditional on a rule fixed in advance. It required, on validation data and ImageNet features, (1) a band with $dS\ge 0.05$ (lower interval bound above zero) and $dD\le0.02$ (upper bound below 0.03); (2) a matched-energy random band-stop control with a smaller $dS$; (3) in an exploratory C2-lite step (a linear head on band-stopped features, no training), a source-accuracy drop of at least 0.05 with a within-source macro-F1 loss of at most 0.02; and (4) the student's approval of the compute. Otherwise the decision is NO-GO and C2 is listed as future work. C2-lite used the band with the largest $dS$ and a larger training sample.

\subsection{Evaluation order, access log and wording rules}
\label{sec:eval-order}
The decision rules, thresholds, variant list, sample sizes and seeds of C1--C4 were written down and approved before any novelty result existed. They were committed to the repository and any later change would have to be logged with its reason; no change was made to a setting after a result existed. ``Fixed in advance and recorded in the repository'' is the strongest statement we make; they were not registered with any external body. Each analysis computed its test and Montgomery metrics once, in a logged order (C1 main, C1 trained backbone, C1 no-filter, C3, C4, C2-lite), and the log lists each access with its date, code commit and output file. Montgomery was scored exactly twice by this extension (C3 and C4), none in C1 or C2-lite. With 13 variants, six bands, five folds and up to 20 iterations, many intervals are computed; no p-values are computed, and a statement about one band, fold or iteration is made only if its interval excludes the reference and the same pattern holds across neighbouring bands or the secondary extractor.
